\documentclass[12pt,a4paper]{article}

% --- Pacotes Básicos ---
\usepackage[utf8]{inputenc}
\usepackage[portuguese]{babel}
\usepackage[left=3cm,right=2cm,top=3cm,bottom=2cm]{geometry}
\usepackage{booktabs} % Para tabelas bonitas
\usepackage{tabularx} % Tabelas com largura ajustada
\usepackage{lipsum}   % Para gerar texto de preenchimento (encher as 5 páginas)
\usepackage{xcolor}   % Para cores nos status
\usepackage{fancyhdr} % Cabeçalhos e rodapés

% --- Configuração de Cabeçalho ---
\pagestyle{fancy}
\fancyhf{}
\lhead{Documentação Técnica Médica}
\rhead{Controle de Suprimentos}
\cfoot{\thepage}

\title{\textbf{Relatório Técnico de Gestão de Insumos Médicos}\\ \large Controle de Estoque, Recebimento e Validade}
\author{Departamento de Logística Hospitalar}
\date{25 de Março de 2026}

\begin{document}

\maketitle
\newpage

\tableofcontents
\newpage

% --- PÁGINA 1: INTRODUÇÃO ---
\section{Introdução e Objetivos}
Este documento estabelece as diretrizes técnicas para o manejo de suprimentos e medicamentos. A rastreabilidade, desde a data de chegada até o consumo final, é o pilar para garantir a segurança do paciente.

\subsection{Importância do Controle de Vencimento}
O descarte de medicamentos por validade expirada representa uma perda financeira significativa e um risco sanitário. Implementamos aqui o sistema \textbf{PVPS} (Primeiro que Vence, Primeiro que Sai).

\section{Metodologia de Recebimento}
Todo insumo que chega à unidade deve passar por conferência tripla:
\begin{enumerate}
    \item Integridade da embalagem.
    \item Concordância entre Nota Fiscal e Quantidade Física.
    \item Registro imediato do lote e validade no sistema central.
\end{enumerate}

\newpage

% --- PÁGINA 2: TABELAS DE SUPRIMENTOS ---
\section{Inventário Detalhado: Medicamentos}
Abaixo, apresenta-se a relação atualizada dos principais fármacos em estoque, com foco nas datas críticas de reposição e expiração.

\begin{table}[h]
\centering
\begin{tabularx}{\textwidth}{|X|c|c|c|c|}
\hline
\textbf{Medicamento} & \textbf{Qtd.} & \textbf{Chegada} & \textbf{Vencimento} & \textbf{Status} \\ \hline
Insulina Regular & 500 un & 10/01/2026 & 15/05/2027 & \textcolor{green}{Seguro} \\ \hline
Amoxicilina 500mg & 1200 cap & 20/02/2026 & 10/11/2026 & \textcolor{orange}{Alerta} \\ \hline
Morfina 10mg/ml & 50 amp & 05/03/2026 & 20/01/2028 & \textcolor{green}{Seguro} \\ \hline
Dipirona Sódica & 3000 ml & 15/03/2026 & 14/04/2026 & \textcolor{red}{Crítico} \\ \hline
Soro Fisiológico & 100 bags & 18/03/2026 & 01/12/2028 & \textcolor{green}{Seguro} \\ \hline
\end{tabularx}
\caption{Controle de Estoque e Validade - Setor A1}
\end{table}

\section{Suprimentos Hospitalares (Descartáveis)}
Diferente dos remédios, os suprimentos de consumo (como luvas e seringas) possuem maior rotatividade.

\begin{itemize}
    \item \textbf{Seringas 5ml:} 5.000 unidades (Chegada: 01/03/2026).
    \item \textbf{Luvas de Procedimento:} 200 caixas (Vencimento Médio: 2029).
    \item \textbf{Cateteres IV:} 300 unidades (Data de Chegada: 12/03/2026).
\end{itemize}

\newpage

% --- PÁGINA 3: PROTOCOLOS ---
\section{Protocolo de Armazenagem e Temperatura}
Muitos dos medicamentos listados na página anterior exigem controle térmico rigoroso.

\subsection{Cadeia de Frio}
Os itens com vencimento próximo devem ser movidos para a frente das prateleiras de refrigeração. A temperatura deve ser mantida entre $2^\circ C$ e $8^\circ C$.

\subsection{Fluxo de Reposição}
A reposição ocorre sempre que o estoque atinge 20\% da capacidade máxima. 
\lipsum[1-2] % Preenchendo com texto técnico simulado

\newpage

% --- PÁGINA 4: ANÁLISE DE RISCO ---
\section{Análise de Riscos e Perdas}
Nesta seção, detalhamos o que acontece quando um medicamento atinge a data de vencimento sem ser utilizado.

\subsection{Segregação de Itens Vencidos}
Itens vencidos devem ser colocados em caixas vermelhas e lacrados para incineração, conforme a norma RDC 222/2018 da ANVISA.

\subsection{Relatório de Desperdício}
\lipsum[3-5] % Mais texto para garantir as 5 páginas

\newpage

% --- PÁGINA 5: CONCLUSÃO ---
\section{Conclusão e Próximos Passos}
A manutenção desta documentação técnica é vital para a auditoria hospitalar. 

\subsection{Recomendações}
\begin{itemize}
    \item Digitalizar todos os canhotos de chegada.
    \item Realizar contagem física (inventário rotativo) a cada 15 dias.
    \item Revisar os fornecedores com prazos de entrega superiores a 10 dias.
\end{itemize}

\vspace{2cm}
\begin{center}
\rule{8cm}{0.4pt} \\
Assinatura do Responsável Técnico \\
CRF/CRM Ativo
\end{center}

\end{document}